% questions.tex —— 1954 年·全国卷 · 题目内容片段（无答案）
% 由 tools/generate_projects.py 依来源总账自 1954/1954_数学真题_试题.md 生成，勿手改（改源 MD 后重跑）


\par\medskip

\gkpart{一、小题}

\begin{question}
\gkti{1}（甲）化简 $\left[\left(a^{\frac32}\cdot b^{-2}\right)\cdot\left(a\cdot b^{-3}\right)^{\frac12}\cdot\left(b^{\frac12}\right)^7\right]^{\frac13}$．

\end{question}
\begin{question}
\gkti{2}（乙）解 $\dfrac16\log x=\dfrac13\log a+2\log b+\log c$．

\end{question}
\begin{question}
\gkti{3}（丙）用二项式定理计算 $(3.02)^4$，使误差小于 $\dfrac1{1000}$．

\end{question}
\begin{question}
\gkti{4}（丁）直角三角形弦上半圆的面积等于勾上半圆与股上半圆面积之和，试证明之．

\end{question}
\begin{question}
\gkti{5}（戊）已知球的半径为 $r$，求内接正方体的体积．

\end{question}
\begin{question}
\gkti{6}（己）已知三角形的一边之长为 $a$，两邻角为 $\beta$ 及 $\gamma$，求计算边长 $b$ 的计算公式．

\par\medskip

\end{question}
\gkpart{二、大题}

\begin{question}
\gkti{7}（二）描绘 $y=3x^2-7x-1$ 之图象，并按下列条件分别求 $x$ 的值的范围：（i）$y>0$；（ii）$y<0$．

\end{question}
\begin{question}
\gkti{8}（三）假设两圆互相外切，求证用连心线段为直径所作的圆必与前两圆的外公切线相切．

\end{question}
\begin{question}
\gkti{9}（四）解 $\dfrac{1+\tan x}{1-\tan x}=1+\sin2x$，求 $x$ 之通值．

\end{question}
\begin{question}
\gkti{10}（五）有一直圆锥，全面积为 $a$；与之同底同高之直圆柱全面积为 $a'$．求该圆锥高与母线之比．

\end{question}